% 原分章PDF93页附录7.D标题—94页末，印刷页669—670。
\section*{附录7.D\quad 反向的M–L距离}
\addcontentsline{toc}{section}{附录7.D 反向的M–L距离}
\markright{附录7.D 反向的M–L距离}

利用式（7.66），得到
\begin{equation}
D(P_X\parallel U_X)=D(P_X\parallel V_{out}(X))
 -D(U_X\parallel V_{out}(X)).\tag{7.D.1}
\end{equation}
注意
\begin{equation}
-D(U_X\parallel V_{out})<D(V_{out}(X)\parallel U_X).\tag{7.D.2}
\end{equation}
将式（7.D.2）代入式（7.D.1），得到
\begin{equation}
D(P_X\parallel V_{out}(X))+D(V_{out}(X)\parallel U_X)
 >D(P_X\parallel U_X).\tag{7.D.3}
\end{equation}
回顾
\begin{equation}
V_{out}(x)=V_Yf'(x),\qquad U_{out}(Y)=\frac{U_X}{f'(x)}.\tag{7.D.4}
\end{equation}
利用式（7.D.4），得到
\begin{equation}
\begin{aligned}
D(V_Y\parallel U_{out}(Y))
 &=\int\log\frac{V_Y}{U_{out}(Y)}V_Y\,dy\\
 &=\int\log\frac{V_Yf'(x)}{U_X}V_Yf'(x)\,dx\\
 &=\int\log\frac{V_{out}(X)}{U_X}V_{out}(X)\,dx\\
 &=D(V_{out}(X)\parallel U_X).
\end{aligned}\tag{7.D.5}
\end{equation}
利用式（7.C.7）、（7.D.5），得到
\begin{equation}
\begin{aligned}
D(P_Y\parallel V_Y)+D(V_Y\parallel U_{out}(Y))
 ={}&D(P_X\parallel V_{out}(X))\\
&{}+D(V_{out}(X)\parallel U_X).
\end{aligned}\tag{7.D.6}
\end{equation}
由式（7.C.8）、（7.D.6）及式（7.70），得到
\begin{equation}
D(P_Y\parallel V_Y)+D(V_Y\parallel U_{out}(Y))
 >D(P_Y\parallel U_{out}(Y)),\tag{7.D.7}
\end{equation}
从而证明了不等式（7.71）。

现在计算 $\operatorname{cov}_{V_Y}\bigl(\log(U_{out}(Y)/V_Y),P_Y/V_Y\bigr)$。由协方差的定义，有
\begin{equation}
\begin{aligned}
&\operatorname{cov}_{V_Y}\left(\log\frac{U_{out}(Y)}{V_Y},\frac{P_Y}{V_Y}\right)\\
&\quad=\int\log\frac{U_{out}(Y)}{V_Y}\frac{P_Y}{V_Y}V_Y\,dy\\
&\qquad\quad{}-\int\log\frac{U_{out}(Y)}{V_Y}V_Y\,dy
 \int\frac{P_Y}{V_Y}V_Y\,dy\\
&\quad=\int\log\frac{U_{out}(Y)}{V_Y}P_Y\,dy
 -\int\log\frac{U_{out}(Y)}{V_Y}V_Y\,dy\\
&\quad=\int\log\frac{U_{out}(Y)}{P_Y}\frac{P_Y}{V_Y}P_Y\,dy
 +D(V_Y\parallel U_{out}(Y))\\
&\quad=D(P_Y\parallel V_Y)-D(P_Y\parallel U_{out}(Y))
 +D(V_Y\parallel U_{out}(Y)).
\end{aligned}\tag{7.D.8}
\end{equation}
利用不等式（7.D.7），式（7.D.8）表明
\[
\operatorname{cov}_{V_Y}\left(\log\frac{U_{out}(Y)}{V_Y},\frac{P_Y}{V_Y}\right)>0.
\]
结论7.10中的其他不等式可类似证明。
